% Generated draft: requires mathematical and source review.
\documentclass[11pt]{article}
\usepackage[margin=25mm]{geometry}
\usepackage{fontspec}
\setmainfont{Latin Modern Roman}
\usepackage{amsmath,amssymb,mathtools,unicode-math}
\setmathfont{Latin Modern Math}
\usepackage{enumitem,needspace,xurl,hyperref}
\setlength{\parindent}{0pt}
\setlength{\parskip}{5pt}
\setlength{\emergencystretch}{3em}
\newcommand{\N}{\mathbb{N}}
\newcommand{\Z}{\mathbb{Z}}
\newcommand{\Q}{\mathbb{Q}}
\newcommand{\R}{\mathbb{R}}
\newcommand{\C}{\mathbb{C}}
\newcommand{\F}{\mathbb{F}}
\newcommand{\A}{\mathbb{A}}
\newcommand{\PP}{\mathbb{P}}
\newcommand{\E}{\mathbb{E}}
\newcommand{\eps}{\varepsilon}
\newcommand{\dd}{\,\mathrm{d}}
\begin{document}
\section*{UnsolvedMath ID: 30005259}
\textbf{Source problem code:} catalogue entry
\subsection{Definitions and mathematical statement}
Let \((S,d)\) be a metric space, and let \(\mathcal B(S)\) denote its Borel \(\sigma\)-algebra. A metric measure space additionally carries a measure \(\mu\); the supplied record does not specify its normalization or its relationship to the sampling law. For observations \(X_1,\ldots,X_n\) taking values in \(S\), define their empirical probability measure by
\[
P_n=\frac1n\sum_{i=1}^n\delta_{X_i},
\]
where \(\delta_x\) is the unit point mass at \(x\).

\textbf{Problem as supplied.} For object-valued observations in a metric measure space, specify an appropriate population mean and its empirical counterpart, and determine hypotheses under which the empirical mean satisfies a law of large numbers and a central limit theorem.

For example, \emph{if} the intended construction assigns a unique point \(m(P)\in S\) to a probability law \(P\), and \(m(P_n)\) is defined, the consistency question can be expressed schematically as
\[
d\bigl(m(P_n),m(P)\bigr)\longrightarrow 0
\qquad(n\longrightarrow\infty),
\]
with the mode of convergence and sampling assumptions still to be specified. This is a conditional template, not an identification of the source's intended mean or a claimed theorem.

\textbf{Unresolved qualifications.} The source extract does not identify the antecedent of ``it,'' define a mean, or specify whether a mean is unique or set-valued. It supplies no independence, identical-distribution, moment, measurability, geometric, or selection hypotheses. It also does not specify the normalization, representation of fluctuations, limiting distribution, or topology needed for a central limit theorem. In a general metric space, subtraction of mean points has no prescribed meaning. Consequently, the supplied evidence supports this research question but does not determine a unique, fully quantified limit-theorem conjecture.
\subsection{Short English statement}
The problem asks how to define a mean for data taking values in a metric space, and when empirical means obey a law of large numbers and a central limit theorem. The supplied extract leaves the intended mean, sampling assumptions, convergence mode, and fluctuation framework unspecified, so a precise theorem requires further source clarification.
\subsection{Sources}
\begin{itemize}
\item Local registry: \texttt{data/unsolvedmath/problems.json}, permanent UnsolvedMath ID \(30005259\), title ``Limit Theorems for Object Data on Metric Spaces,'' source problem code \texttt{catalogue entry}. This draft uses the supplied complete registry record.
\item Supplied report citation: \emph{Heat Kernels, Stochastic Processes and Functional Inequalities} (2023), Oberwolfach Reports, pp.\ 2811--2893, DOI \texttt{10.4171/owr/2022/49}. Supplied source URL: \url{https://doi.org/10.4171/owr/2022/49}. The record assigns proposed year 2022; the relationship between that year and the citation's 2023 date is not explained.
\item Supplied literature source: ``Oberwolfach Report 49,'' \url{https://ems.press/content/serial-article-files/46987}. Literature assessment checked at \texttt{2026-08-22}: the supplied curation labels the problem partially solved and reports partial answers to the law of large numbers and central limit questions, but no general theorem settling the extracted problem. This is supplied dated curation, not an independently verified assessment.
\item Attribution in the supplied original extract: ``Object data analysis -- statistics on metric spaces,'' Thomas Hotz; followed by ``Gromov--Wasserstein Distances, Distance Distributions and Inverse Problems for Metric Measure Spaces,'' Tom Needham, joint work with Facundo M\'{e}moli. The extract does not establish whether the latter heading belongs to the question. The registry supplies no named proposer.
\item The registry labels the statement \texttt{corrected\_verified} and says that the official report was checked, the question isolated, and notation repaired. These are supplied verification claims; no independent source consultation is asserted here.
\end{itemize}
\end{document}
